\documentclass[a4paper,12pt]{article}
\usepackage[utf8]{inputenc}
\usepackage[ngerman]{babel}
\usepackage{amsmath}
\usepackage{parskip}

\setlength{\parindent}{0pt}

\title{Modul 293 - Aufgabenblatt 11}
\author{von Lukas Meier}
\date{Unterricht vom 17.03.2026}

\begin{document}

\maketitle

\section*{Aufgabenblatt: HTML-Formulare I}

\subsection*{Teil 1: Grundbegriffe verstehen}
Beantworte die folgenden Fragen kurz:
\begin{itemize}
  \item Wofür wird das Element \texttt{form} verwendet?
  \item Was macht das Attribut \texttt{action} ganz grob?
  \item Was macht das Attribut \texttt{method} ganz grob?
  \item Welche Aufgabe hat ein Button mit \texttt{type="submit"}?
\end{itemize}

\subsection*{Teil 2: Formular-Grundstruktur schreiben}
Erstelle ein einfaches HTML-Formular mit:
\begin{itemize}
  \item einem \texttt{form}-Element
  \item \texttt{action="/anfrage"}
  \item \texttt{method="post"}
  \item einem Submit-Button
\end{itemize}

\subsection*{Teil 3: Eingabefelder ergänzen}
Erweitere dein Formular aus Teil 2 mit:
\begin{itemize}
  \item einem Textfeld für den Vornamen
  \item einem Textfeld für den Nachnamen
  \item einem Zahlenfeld für das Alter
\end{itemize}

Verwende für jedes Feld sinnvolle \texttt{name}-Werte.

\subsection*{Teil 4: Code lesen}
Betrachte den folgenden Code:

\begin{verbatim}
<form action="/send" method="get">
  <input type="text" name="stadt">
  <input type="number" name="plz">
  <button>Senden</button>
</form>
\end{verbatim}

\begin{itemize}
  \item Welche Daten kann der Benutzer hier eingeben?
  \item Was passiert beim Klick auf den Button?
\end{itemize}

\subsection*{Teil 5: Fehler korrigieren}
Der folgende Code enthält mehrere Probleme:

\begin{verbatim}
<form action="/registrieren" method="post">
  <input type="text" name="vorname"
  <input type="number" name="alter"></input>
  <button type="submit">Absenden
</form>
\end{verbatim}

\begin{itemize}
  \item Finde mindestens drei Fehler.
  \item Korrigiere den Code vollständig.
  \item Begründe kurz jede Korrektur.
\end{itemize}

\subsection*{Teil 6: Praxisaufgabe}
Erstelle eine kleine HTML-Seite mit einem Formular für eine Kursanmeldung.

Anforderungen:
\begin{itemize}
  \item Überschrift und kurzer Einleitungstext
  \item ein Formular mit \texttt{action} und \texttt{method}
  \item zwei Textfelder
  \item ein Zahlenfeld
  \item ein Button mit \texttt{type="submit"}
\end{itemize}

Thema-Vorschlag: Anmeldung für einen HTML-Workshop.

\subsection*{Teil 7: Reflexion}
Beantworte kurz:
\begin{itemize}
  \item Was ist der Unterschied zwischen dem Formular selbst und den Eingabefeldern?
  \item Warum braucht ein Formular meistens einen Submit-Button?
  \item Was war heute noch unklar?
\end{itemize}

\end{document}
